% Shared IEEE conference formatting. Keep the class, margins and body size intact.
\documentclass[conference,letterpaper]{IEEEtran}
\usepackage[T1]{fontenc}
\usepackage{amsmath,amssymb}
\usepackage{cite}
\usepackage{booktabs,array}
\usepackage{balance}
\usepackage[hidelinks]{hyperref}
\urlstyle{same}
\hypersetup{pdfdisplaydoctitle=true}

\hypersetup{
  pdftitle={Muscle Mechanics for Aided Diagnosis},
  pdfauthor={Santiago Restrepo Ruiz},
  pdfsubject={Regional biomechanics and guided movement observation},
  pdfkeywords={biomechanics, musculoskeletal modeling, static optimization,
    WebAssembly, movement observation}
}

\title{Muscle Mechanics for Aided Diagnosis}
\author{\IEEEauthorblockN{Santiago Restrepo Ruiz}
\IEEEauthorblockA{Published by Waajacu Open Source Foundation\\
\href{mailto:Santiago.Restrepo.Ruiz@gmail.com}{Santiago.Restrepo.Ruiz@gmail.com}\quad
\url{https://medical.waajacu.com/}\\[3pt]
\textit{Regional biomechanics and guided movement observation}}}
\date{}

\begin{document}
\maketitle

\begin{abstract}
We developed Muscles, a browser application that connects interactive anatomy,
regional static muscle estimates and structured movement observations. It
adapts MS-Human-700 to a client-side MuJoCo runtime, recomputes muscle paths for
each posture, and displays activation estimates only when explicit mechanical
checks pass. Versioned observation protocols and local reports support
reproducible exploration. The system is intended for research and education;
its generic model outputs are not patient measurements or diagnoses.
\end{abstract}

\begin{IEEEkeywords}
Biomechanics, musculoskeletal modeling, static optimization.
\end{IEEEkeywords}

\balance
\section{Browser Architecture and Model Integration}
The application builds on MS-Human-700~\cite{zuo2024self} and
MuJoCo~\cite{mujoco}. Our contribution is the browser integration, regional
calculation workflow, interactive visualization and observation/reporting
layer. Pinned model files are compiled into browser assets with provenance
and integrity hashes. Two documented bilateral path-coordinate corrections
are maintained in the imported model.

MuJoCo WebAssembly runs in a Web Worker, while Three.js renders the skeleton
and muscle paths. The server delivers static files; model computation runs
locally. Users select either arm, either leg, either foot and ankle, trunk,
head and neck, or the detailed right hand. Each selection supplies its own
coordinates, muscles and posture presets. The primary model supplies eight
regions; a separate manipulation profile supplies the articulated right hand.
The selected region is emphasized within a complete skeleton for context.

\section{Static Muscle Estimation}
For each posture, the system resets the model, applies the selected joint
coordinates and their equality dependents, and updates tendon paths with
their via points, wrapping contacts and pulley segments. Other coordinates
remain prescribed. At zero velocity and acceleration, reduced-coordinate
force balance includes gravity, passive forces and muscle actuation.

Let \(a\) denote muscle activation, \(A(q)\) the activation-to-torque map at
posture \(q\), \(b(q)\) the zero-activation torque baseline, and \(r\) unbounded reserve
torque. The implemented bounded optimization is
\begin{equation*}
\begin{aligned}
\min_{a,r}\quad & \lVert a\rVert_2^2+\lambda\lVert r\rVert_2^2\\
\text{subject to}\quad & A(q)a+r=-b(q),\qquad 0\leq a\leq1.
\end{aligned}
\end{equation*}
The solver uses \(\lambda=2500\,(\mathrm{N\,m})^{-2}\) and replays the
solution through MuJoCo. Activation colors require a maximum absolute residual
no greater than \(10^{-4}\,\mathrm{N\,m}\), reserve torque no greater than
\(0.05\,\mathrm{N\,m}\), and passing path, finite-value and capacity checks.
Failed checks suppress activation estimates. Reserve use indicates unmet
modeled torque demand; it does not establish physiological weakness.

\section{Guided Observation and Reports}
Guided Assessment selects a body region and, where available, its explicit
side, then loads the matching versioned protocol. The right-arm protocol
contains 20 distinct postures selected from a deterministic 384-candidate
atlas; the final five are supplemental. Selection sought posture variety and
distinct generic activation patterns among mechanically acceptable candidates.
Eight additional regional protocols extend the observation workflow.

Users record movement completion and symptoms separately from model demand.
Local, versioned records support saved reports and matched comparisons;
protocol changes archive incompatible observations and exclude them from the
active protocol. The main export removes direct identifiers and free text. Six interface
languages share stable protocol identifiers and report contracts.

\section{Verification and Scope}
Repository verification checks model hashes, regional inventories, protocol
identity, localization consistency, deployment files and HTTP behavior.
The report suite passes 20 tests, and protocol checks cover nine definitions.
These checks establish implementation consistency, not clinical validity.

The calculation describes an unloaded static hold with prescribed support.
It does not estimate gait, ground reaction, grip force, fatigue, injury or
patient-specific tissue loading. Generic activation is not measured effort,
and recorded symptoms do not identify a painful muscle. The application is
research and educational software, not a medical device.

Explore the application and its available tool routes at
\mbox{\url{https://medical.waajacu.com/}}.

\bibliographystyle{IEEEtran}
\bibliography{references}
\end{document}
